\documentclass{amsart}
% For dealing with references we use the comment environment
\usepackage{verbatim}
\newenvironment{reference}{\comment}{\endcomment}
%\newenvironment{reference}{}{}
\newenvironment{slogan}{\comment}{\endcomment}
\newenvironment{history}{\comment}{\endcomment}

% For commutative diagrams we use Xy-pic
\usepackage[all]{xy}

% We use 2cell for 2-commutative diagrams.
\xyoption{2cell}
\UseAllTwocells

% We use multicol for the list of chapters between chapters
\usepackage{multicol}

% This is generally recommended for better output
\usepackage{lmodern}
\usepackage[T1]{fontenc}

% For cross-file-references

% Package for hypertext links:
\usepackage{hyperref}

% For any local file, say "hello.tex" you want to link to please

% Theorem environments.
%
\theoremstyle{plain}
\newtheorem{theorem}[subsection]{Theorem}
\newtheorem{proposition}[subsection]{Proposition}
\newtheorem{lemma}[subsection]{Lemma}

\theoremstyle{definition}
\newtheorem{definition}[subsection]{Definition}
\newtheorem{example}[subsection]{Example}
\newtheorem{exercise}[subsection]{Exercise}
\newtheorem{situation}[subsection]{Situation}

\theoremstyle{remark}
\newtheorem{remark}[subsection]{Remark}
\newtheorem{remarks}[subsection]{Remarks}

\numberwithin{equation}{subsection}

% Macros
%
\def\lim{\mathop{\mathrm{lim}}\nolimits}
\def\colim{\mathop{\mathrm{colim}}\nolimits}
\def\Spec{\mathop{\mathrm{Spec}}}
\def\Hom{\mathop{\mathrm{Hom}}\nolimits}
\def\Ext{\mathop{\mathrm{Ext}}\nolimits}
\def\SheafHom{\mathop{\mathcal{H}\!\mathit{om}}\nolimits}
\def\SheafExt{\mathop{\mathcal{E}\!\mathit{xt}}\nolimits}
\def\Sch{\mathit{Sch}}
\def\Mor{\mathop{\mathrm{Mor}}\nolimits}
\def\Ob{\mathop{\mathrm{Ob}}\nolimits}
\def\Sh{\mathop{\mathit{Sh}}\nolimits}
\def\NL{\mathop{N\!L}\nolimits}
\def\CH{\mathop{\mathrm{CH}}\nolimits}
\def\proetale{{pro\text{-}\acute{e}tale}}
\def\etale{{\acute{e}tale}}
\def\QCoh{\mathit{QCoh}}
\def\Ker{\mathop{\mathrm{Ker}}}
\def\Im{\mathop{\mathrm{Im}}}
\def\Coker{\mathop{\mathrm{Coker}}}
\def\Coim{\mathop{\mathrm{Coim}}}

% Boxtimes
%
\DeclareMathSymbol{\boxtimes}{\mathbin}{AMSa}{"02}

%
% Macros for moduli stacks/spaces
%
\def\QCohstack{\mathcal{QC}\!\mathit{oh}}
\def\Cohstack{\mathcal{C}\!\mathit{oh}}
\def\Spacesstack{\mathcal{S}\!\mathit{paces}}
\def\Quotfunctor{\mathrm{Quot}}
\def\Hilbfunctor{\mathrm{Hilb}}
\def\Curvesstack{\mathcal{C}\!\mathit{urves}}
\def\Polarizedstack{\mathcal{P}\!\mathit{olarized}}
\def\Complexesstack{\mathcal{C}\!\mathit{omplexes}}
% \Pic is the operator that assigns to X its picard group, usage \Pic(X)
% \Picardstack_{X/B} denotes the Picard stack of X over B
% \Picardfunctor_{X/B} denotes the Picard functor of X over B
\def\Pic{\mathop{\mathrm{Pic}}\nolimits}
\def\Picardstack{\mathcal{P}\!\mathit{ic}}
\def\Picardfunctor{\mathrm{Pic}}
\def\Deformationcategory{\mathcal{D}\!\mathit{ef}}

\setlength{\emergencystretch}{3em}
\begin{document}
\title[Stacks Project, Tag 07CR]{731 --- Strictly standard Noetherian algebras with compatible mod-four sections admit desingularizing factorizations}
\maketitle
\noindent Copyright (C) 2005--2025 Johan de Jong. Stacks Project authors. Source: \href{https://stacks.math.columbia.edu/tag/07CR}{Tag 07CR}.

\noindent Editorial difficulty note (not author text). Difficulty H3: The proof constructs auxiliary equations from Jacobian minors and a mod-four section, then controls residual equations using Krull intersection and annihilator quotients. The explicit factorization simultaneously improves the smooth locus and preserves a map to the target algebra..

\noindent Editorial provenance note (not author text). History: excerpt from revision \texttt{a0444\allowbreak 6e57e\allowbreak c1fbc\allowbreak 252a8\allowbreak 71afc\allowbreak ec775\allowbreak 2fb28\allowbreak 07b14}, smoothing.tex, lines 1664--1809. Reference presentation was adapted; original-author context and core support blocks were retained or added. The mathematical wording is unchanged. Licensed under GNU FDL 1.2 or later; no invariant sections or cover texts. See ../licenses/COPYING. Context blocks reproduce author definitions and supporting results; they are not additional counted problems.

\begin{definition}
\label{definition-strictly-standard}
Let $R \to A$ be a ring map of finite presentation.
We say an element $a \in A$ is {\it elementary standard in $A$ over $R$}
if there exists a presentation
$A = R[x_1, \ldots, x_n]/(f_1, \ldots, f_m)$
and $0 \leq c \leq \min(n, m)$ such that
\begin{equation}
\label{equation-elementary-standard-one}
a = a' \det(\partial f_j/\partial x_i)_{i, j = 1, \ldots, c}
\end{equation}
for some $a' \in A$ and
\begin{equation}
\label{equation-elementary-standard-two}
a f_{c + j} \in (f_1, \ldots, f_c) + (f_1, \ldots, f_m)^2
\end{equation}
for $j = 1, \ldots, m - c$. We say $a \in A$ is
{\it strictly standard in $A$ over $R$} if there exists a presentation
$A = R[x_1, \ldots, x_n]/(f_1, \ldots, f_m)$
and $0 \leq c \leq \min(n, m)$ such that
\begin{equation}
\label{equation-strictly-standard-one}
a = \sum\nolimits_{I \subset \{1, \ldots, n\},\ |I| = c}
a_I \det(\partial f_j/\partial x_i)_{j = 1, \ldots, c,\ i \in I}
\end{equation}
for some $a_I \in A$ and
\begin{equation}
\label{equation-strictly-standard-two}
a f_{c + j} \in (f_1, \ldots, f_c) + (f_1, \ldots, f_m)^2
\end{equation}
for $j = 1, \ldots, m - c$.
\end{definition}

\begin{definition}
\label{definition-singular-ideal}
Let $R \to A$ be a ring map. The {\it singular ideal of $A$ over $R$},
denoted $H_{A/R}$ is the unique radical ideal $H_{A/R} \subset A$ with
$$
V(H_{A/R}) = \{\mathfrak q \in \Spec(A) \mid R \to A
\text{ not smooth at }\mathfrak q\}
$$
\end{definition}

\begin{lemma}[Elkik]
\label{lemma-elkik}
Let $R \to A$ be a ring map of finite presentation.
The singular ideal $H_{A/R}$ is the radical of the ideal
generated by strictly standard elements in $A$ over $R$
and also the radical of the ideal generated by elementary
standard elements in $A$ over $R$.
\end{lemma}

\begin{proof}
Assume $a$ is strictly standard in $A$ over $R$. We claim that
$A_a$ is smooth over $R$, which proves that $a \in H_{A/R}$. Namely,
let $A = R[x_1, \ldots, x_n]/(f_1, \ldots, f_m)$, $c$, and $a' \in A$
be as in Definition \ref{definition-strictly-standard}.
Write $I = (f_1, \ldots, f_m)$ so that the naive cotangent
complex of $A$ over $R$ is given by $I/I^2 \to \bigoplus A\text{d}x_i$.
Assumption (\ref{equation-strictly-standard-two})
implies that $(I/I^2)_a$ is generated by the classes of $f_1, \ldots, f_c$.
Assumption (\ref{equation-strictly-standard-one}) implies
that the differential $(I/I^2)_a \to \bigoplus A_a\text{d}x_i$
has a left inverse, see
Lemma \href{https://stacks.math.columbia.edu/tag/07ET}{lemma parse equation strictly standard one (Tag 07ET)}.
Hence $R \to A_a$ is smooth by definition and
Algebra, Lemma \href{https://stacks.math.columbia.edu/tag/00S7}{lemma localize NL (Tag 00S7)}.

\medskip\noindent
Let $H_e, H_s \subset A$ be the radical of the ideal generated by
elementary, resp.\ strictly standard elements of $A$ over $R$.
By definition and what we just proved we have
$H_e \subset H_s \subset H_{A/R}$. The inclusion $H_{A/R} \subset H_e$
follows from Lemma \href{https://stacks.math.columbia.edu/tag/07C6}{lemma find strictly standard (Tag 07C6)}.
\end{proof}


\begin{lemma}
\label{lemma-desingularize}
Let $R$ be a Noetherian ring.
Let $\Lambda$ be an $R$-algebra. Let $\pi \in R$ and
assume that $\text{Ann}_\Lambda(\pi) = \text{Ann}_\Lambda(\pi^2)$. Let
$A \to \Lambda$ be an $R$-algebra map with $A$ of finite
presentation. Assume
\begin{enumerate}
\item the image of $\pi$ is strictly standard in $A$ over $R$, and
\item there exists a section $\rho : A/\pi^4 A \to R/\pi^4 R$
which is compatible with the map to $\Lambda/\pi^4 \Lambda$.
\end{enumerate}
Then we can find $R$-algebra maps $A \to B \to \Lambda$ with $B$
of finite presentation such that $\mathfrak a B \subset H_{B/R}$ where
$\mathfrak a = \text{Ann}_R(\text{Ann}_R(\pi^2)/\text{Ann}_R(\pi))$.
\end{lemma}

\begin{proof}
Choose a presentation
$$
A = R[x_1, \ldots, x_n]/(f_1, \ldots, f_m)
$$
and $0 \leq c \leq \min(n, m)$ such that
(\ref{equation-strictly-standard-one}) holds for $\pi$ and such that
\begin{equation}
\label{equation-star}
\pi f_{c + j} \in (f_1, \ldots, f_c) + (f_1, \ldots, f_m)^2
\end{equation}
for $j = 1, \ldots, m - c$. Say $\rho$ maps $x_i$ to the class of
$r_i \in R$. Then we can replace $x_i$ by $x_i - r_i$. Hence we may
assume $\rho(x_i) = 0$ in $R/\pi^4 R$. This implies that
$f_j(0) \in \pi^4R$ and that $A \to \Lambda$ maps $x_i$
to $\pi^4\lambda_i$ for some $\lambda_i \in \Lambda$. Write
$$
f_j = f_j(0) + \sum\nolimits_{i = 1, \ldots, n} r_{ji} x_i + \text{h.o.t.}
$$
This implies that the constant term of $\partial f_j/\partial x_i$ is
$r_{ji}$. Apply $\rho$ to (\ref{equation-strictly-standard-one})
for $\pi$ and we see that
$$
\pi = \sum\nolimits_{I \subset \{1, \ldots, n\},\ |I| = c}
r_I \det(r_{ji})_{j = 1, \ldots, c,\ i \in I} \bmod \pi^4R
$$
for some $r_I \in R$. Thus we have
$$
u\pi = \sum\nolimits_{I \subset \{1, \ldots, n\},\ |I| = c}
r_I \det(r_{ji})_{j = 1, \ldots, c,\ i \in I}
$$
for some $u \in 1 + \pi^3R$. By
Algebra, Lemma \href{https://stacks.math.columbia.edu/tag/07DQ}{lemma matrix left inverse (Tag 07DQ)}
this implies there exists a $n \times c$ matrix $(s_{ik})$ such that
$$
u\pi \delta_{jk} = \sum\nolimits_{i = 1, \ldots, n} r_{ji}s_{ik}\quad
\text{for all } j, k = 1, \ldots, c
$$
(Kronecker delta). We introduce auxiliary variables
$v_1, \ldots, v_c, w_1, \ldots, w_n$ and we set
$$
h_i = x_i - \pi^2 \sum\nolimits_{j = 1, \ldots c} s_{ij} v_j - \pi^3 w_i
$$
In the following we will use that
$$
R[x_1, \ldots, x_n, v_1, \ldots, v_c, w_1, \ldots, w_n]/
(h_1, \ldots, h_n) = R[v_1, \ldots, v_c, w_1, \ldots, w_n]
$$
without further mention. In
$R[x_1, \ldots, x_n, v_1, \ldots, v_c, w_1, \ldots, w_n]/
(h_1, \ldots, h_n)$ we have
\begin{align*}
f_j & = f_j(x_1 - h_1, \ldots, x_n - h_n) \\
& =
\pi^2 \sum\nolimits_{k = 1}^c
\left(\sum\nolimits_{i = 1}^n r_{ji} s_{ik}\right) v_k
+
\pi^3 \sum\nolimits_{i = 1}^n r_{ji}w_i \bmod \pi^4 \\
& =
\pi^3 v_j + \pi^3 \sum\nolimits_{i = 1}^n r_{ji}w_i \bmod \pi^4
\end{align*}
for $1 \leq j \leq c$. Hence we can choose elements
$g_j \in R[v_1, \ldots, v_c, w_1, \ldots, w_n]$
such that $g_j = v_j + \sum r_{ji}w_i \bmod \pi$
and such that $f_j = \pi^3 g_j$ in the $R$-algebra
$R[x_1, \ldots, x_n, v_1, \ldots, v_c, w_1, \ldots, w_n]/
(h_1, \ldots, h_n)$. We set
$$
B = R[x_1, \ldots, x_n, v_1, \ldots, v_c, w_1, \ldots, w_n]/
(f_1, \ldots, f_m, h_1, \ldots, h_n, g_1, \ldots, g_c).
$$
The map $A \to B$ is clear. We define $B \to \Lambda$ by mapping
$x_i \to \pi^4\lambda_i$, $v_i \mapsto 0$, and $w_i \mapsto \pi \lambda_i$.
Then it is clear that the elements $f_j$ and $h_i$ are mapped to zero
in $\Lambda$. Moreover, it is clear that $g_i$ is mapped to an element
$t$ of $\pi\Lambda$ such that $\pi^3t = 0$ (as $f_i = \pi^3 g_i$ modulo
the ideal generated by the $h$'s). Hence our assumption that
$\text{Ann}_\Lambda(\pi) = \text{Ann}_\Lambda(\pi^2)$ implies that $t = 0$.
Thus we are done if we can prove the statement about smoothness.

\medskip\noindent
Note that $B_\pi \cong A_\pi[v_1, \ldots, v_c]$ because the equations
$g_i = 0$ are implied by $f_i = 0$. Hence $B_\pi$ is smooth over $R$
as $A_\pi$ is smooth over $R$ by the assumption that $\pi$ is strictly
standard in $A$ over $R$, see
Lemma \ref{lemma-elkik}.

\medskip\noindent
Set $B' = R[v_1, \ldots, v_c, w_1, \ldots, w_n]/(g_1, \ldots, g_c)$.
As $g_i = v_i + \sum r_{ji}w_i \bmod \pi$ we see that
$B'/\pi B' = R/\pi R[w_1, \ldots, w_n]$. Hence
$R \to B'$ is smooth of relative dimension $n$ at every
point of $V(\pi)$ by
Algebra, Lemmas
\href{https://stacks.math.columbia.edu/tag/00ST}{lemma localize relative complete intersection (Tag 00ST)} and
\href{https://stacks.math.columbia.edu/tag/00TF}{lemma flat fibre smooth (Tag 00TF)}
(the first lemma shows it is syntomic at those primes, in particular
flat, whereupon the second lemma shows it is smooth).

\medskip\noindent
Let $\mathfrak q \subset B$ be a prime with $\pi \in \mathfrak q$ and
for some $r \in \mathfrak a$, $r \not \in \mathfrak q$.
Denote $\mathfrak q' = B' \cap \mathfrak q$.
We claim the surjection $B' \to B$ induces an isomorphism of local
rings $(B')_{\mathfrak q'} \to B_\mathfrak q$. This will
conclude the proof of the lemma. Note that $B_\mathfrak q$ is the
quotient of $(B')_{\mathfrak q'}$ by the ideal generated by
$f_{c + j}$, $j = 1, \ldots, m - c$. We observe two things:
first the image of $f_{c + j}$ in $(B')_{\mathfrak q'}$ is
divisible by $\pi^2$ and
second the image of $\pi f_{c + j}$ in $(B')_{\mathfrak q'}$
can be written as $\sum b_{j_1 j_2} f_{c + j_1}f_{c + j_2}$ by
(\ref{equation-star}). Thus we see that the image of each $\pi f_{c + j}$
is contained in the ideal generated by the elements $\pi^2 f_{c + j'}$.
Hence $\pi f_{c + j} = 0$ in $(B')_{\mathfrak q'}$ as this is a
Noetherian local ring, see
Algebra, Lemma \href{https://stacks.math.columbia.edu/tag/00IP}{lemma intersect powers ideal module zero (Tag 00IP)}.
As $R \to (B')_{\mathfrak q'}$ is flat we see that
$$
\left(\text{Ann}_R(\pi^2)/\text{Ann}_R(\pi)\right)
\otimes_R (B')_{\mathfrak q'}
=
\text{Ann}_{(B')_{\mathfrak q'}}(\pi^2)/\text{Ann}_{(B')_{\mathfrak q'}}(\pi)
$$
Because $r \in \mathfrak a$ is invertible in
$(B')_{\mathfrak q'}$ we see that this module is zero.
Hence we see that the image of $f_{c + j}$ is zero in
$(B')_{\mathfrak q'}$ as desired.
\end{proof}
\end{document}
